% ===========================================================================
\section{Emission Factors and Parameters}
\label{sec:factors}
% ===========================================================================

Every factor and parameter used here is an IPCC 2006 default, a 2019 Refinement default
(the wastewater nitrous oxide chain) or a documented regional value,
with its source recorded in the versioned parameter table. Three deserve discussion because
the choice was not automatic.

\begin{table}[htbp]
\centering
\footnotesize
\begin{tabular}{@{}llr@{}}
\toprule
\textbf{Parameter} & \textbf{Applies to} & \textbf{Value} \\
\midrule
Generation rate, urban    & 4A, 4C2 & \SI{0.453}{} kg/cap/day \\
Generation rate, rural    & 4A, 4C2 & \SI{0.283}{} kg/cap/day \\
Uncollected fraction, urban & 4A, 4C2 & 0.30 \\
Uncollected fraction, rural & 4A, 4C2 & 0.70 \\
Burned fraction of uncollected & 4C2 & 0.60 \\
\midrule
Decay constant $k$        & 4A  & 0.065 yr$^{-1}$ \\
$\mathrm{DOC_f}$          & 4A  & 0.50 \\
MCF (site management)     & 4A  & 0.60 \\
$F$ (\CHfour{} in landfill gas) & 4A & 0.50 \\
Oxidation factor          & 4A  & 0 \\
Methane recovery          & 4A  & 0 \\
Back-projection to 1960   & 4A  & 40\% of first panel year \\
\midrule
Dry matter                & 4C2 & 0.70 \\
Total carbon              & 4C2 & \SI{0.2195}{} t C/t wet \\
Fossil carbon             & 4C2 & \SI{0.0757}{} t C/t wet \\
Oxidation factor          & 4C2 & 0.71 \\
\midrule
BOD                       & 4D1 & \SI{40}{} g/cap/day \\
$B_0$                     & 4D1 & \SI{0.6}{} kg \CHfour{}/kg BOD \\
MCF, sewered / septic     & 4D1 & 0.50 \\
MCF, latrine / open       & 4D1 & 0.10 / 0 \\
Industrial-commercial correction & 4D1 & 1.25 \\
\bottomrule
\end{tabular}
\caption{Parameters, with the IPCC volume and table recorded against each in
\texttt{data/waste\_params.csv}. The two zeros in the 4A block are deliberate and argued in
Section~\ref{sec:methods}, not defaults left unset. Full row-level provenance---value,
uncertainty, source and basis for every parameter and factor---is the generated table in
\ref{app:factors}.}
\label{tab:params}
\end{table}

\subsection{Two oxidation factors that must not be reconciled}

This sector carries an oxidation factor of \num{0.71} for open burning and \num{0} for
landfill, while the fossil-fuel sectors of the same inventory carry none at all. That looks
like an inconsistency and is not.

The 2006 Guidelines fold complete oxidation into the fuel combustion emission factors
themselves---Table~1.4 carries an explicit oxidation column of \num{1.0} for every fuel---so
applying the 1996-vintage \num{0.98} on top of a 2006 factor double-discounts. Open burning is
different in kind: it genuinely is incomplete, and \num{0.71} is the 2006 value that says so.
Landfill oxidation is different again, describing methane oxidised while passing through a
cover layer, which unmanaged Pakistani sites largely do not have.

Three different numbers describing three different physical situations. Tests pin all three,
so that a future reader who notices the apparent inconsistency and "corrects" it fails the
build rather than silently changing the sector.

\subsection{Fossil carbon, and why only a third of the burning \COtwo{} is reported}

Total carbon in Pakistani municipal waste is \SI{0.2195}{\tonne} per wet tonne, of which
\SI{0.0757}{\tonne}---roughly a third---is fossil. Only that third produces reportable
\COtwo{}; the rest was fixed from the atmosphere by plants within the year and is reported as
a memo item at \SI{3520}{\giga\gram}.

This convention is easy to apply inconsistently, because the same waste produces \CHfour{} and
\NtwoO{} that \emph{are} counted in full regardless of carbon origin. The distinction is not
about the waste but about the atmosphere: biogenic \COtwo{} returns carbon that was recently
removed, while methane from the same carbon is a different molecule with a different radiative
effect and no such offset.

\subsection{What the parameters are worth}

Three of the parameters above dominate the sector's uncertainty, and each is examined as a
published sensitivity rather than left implicit: the sewer MCF (Section~\ref{sec:discussion}),
the BOD rate, and the spin-up back-projection. The generation and collection rates are not
among them---not because they are well known, but because the partition means they move
emissions between categories more than they move the total.
